% !tp-example Basic usage
\documentclass{article}
\begin{document}
\wcTpOcamlViewerQuestion[open,src=/docs/components/ocaml-playground/examples/double.ml]{
  \begin{description}
    \item[{Title}] Double a number
    \item[{Prompt}] Complete double so it returns twice its argument. Edit the source, then submit to run the tests.
    \item[{Feedback}] /docs/components/ocaml-playground/examples/double.test.ml
    \item[{Solution}] Return x * 2.
  \end{description}
}
\end{document}

% !tp-example Javascript API
\documentclass{article}
\begin{document}
\wcScript[type=module]|
import "/components/ocaml/ocaml.js";
import "/components/ocaml-viewer/ocaml-viewer.js";
import "/components/ocaml-viewer-question/ocaml-viewer-question.js";

const example = document.createDocumentFragment();

const element1 = document.createElement("tp-ocaml-viewer-question");
element1.setAttribute("open", "");
element1.setAttribute("src", "/docs/components/ocaml-playground/examples/double.ml");
element1.append(document.createTextNode("\n  "));

const element2 = document.createElement("dl");
element2.append(document.createTextNode("\n    "));

const element3 = document.createElement("dt");
element3.append(document.createTextNode("Title"));
element2.append(element3);

const element4 = document.createElement("dd");
element4.append(document.createTextNode("Double a number"));
element2.append(element4);
element2.append(document.createTextNode("\n    "));

const element5 = document.createElement("dt");
element5.append(document.createTextNode("Prompt"));
element2.append(element5);

const element6 = document.createElement("dd");
element6.append(document.createTextNode("Complete double so it returns twice its argument. Edit the source, then submit to run the tests."));
element2.append(element6);
element2.append(document.createTextNode("\n    "));

const element7 = document.createElement("dt");
element7.append(document.createTextNode("Feedback"));
element2.append(element7);

const element8 = document.createElement("dd");
element8.append(document.createTextNode("/docs/components/ocaml-playground/examples/double.test.ml"));
element2.append(element8);
element2.append(document.createTextNode("\n    "));

const element9 = document.createElement("dt");
element9.append(document.createTextNode("Solution"));
element2.append(element9);

const element10 = document.createElement("dd");
element10.append(document.createTextNode("Return x * 2."));
element2.append(element10);
element2.append(document.createTextNode("\n  "));
element1.append(element2);
element1.append(document.createTextNode("\n"));
example.append(element1);

document.body.append(example);
|
\end{document}

% !tp-example Passing solution
\documentclass{article}
\begin{document}
\wcTpOcamlViewerQuestion[open,src=/docs/components/ocaml-playground/examples/double-solved.ml,test=/docs/components/ocaml-playground/examples/double.test.ml]{
  \begin{description}
    \item[{Title}] Double a number
    \item[{Prompt}] Complete double so it returns twice its argument. Edit the source, then submit to run the tests.
    \item[{Solution}] Return x * 2.
  \end{description}
}
\end{document}

% !tp-example Explicit Feedback file
\documentclass{article}
\begin{document}
\wcTpOcamlViewerQuestion[open]{
  \wcScript[type=tp/ocaml]|
    let double x = x + 1;;
  |
  \begin{description}
    \item[{Title}] Double a number
    \item[{Prompt}] Complete double so it returns twice its argument. Edit the source, then submit to run the tests.
    \item[{Feedback}] /docs/components/ocaml-playground/examples/double.test.ml
    \item[{Solution}] Return x * 2.
  \end{description}
}
\end{document}
